\chapter{Number theory}

\section{Modular arithmetic}
	\kactlimport{ModularArithmetic.h}
	\kactlimport{ModInverse.h}
	\kactlimport{ModPow.h}
	\kactlimport{ModLog.h}
	\kactlimport{ModSum.h}
	\kactlimport{ModMulLL.h}
	\kactlimport{ModSqrt.h}

\section{Primality}
	\kactlimport{FastEratosthenes.h}
	\kactlimport{MillerRabin.h}
	\kactlimport{Factor.h}

	\textbf{Min\_25 Sieve.}
	Let $f$ be multiplicative, $p_0=1$, $p_k$ the $k$-th prime,
	and $\operatorname{lpf}(m)$ the smallest prime factor of $m$.

	\textbf{Definitions.}
	$F_k(x)=\sum_{2\le m\le x,\ \operatorname{lpf}(m)\ge p_k}f(m)$,
	$F_{\mathrm{prime}}(x)=\sum_{p\le x,\ p\text{ prime}}f(p)$.

	\textbf{Recurrence of $F_k$.}
	\[
	\begin{aligned}
	F_k(x)={}&F_{\mathrm{prime}}(x)-F_{\mathrm{prime}}(p_{k-1})\\
	&+\sum_{\substack{i\ge k\\p_i^2\le x}}
	\sum_{\substack{c\ge1\\p_i^{c+1}\le x}}
	\bigl(f(p_i^c)F_{i+1}(\lfloor x/p_i^c\rfloor)
	+f(p_i^{c+1})\bigr).
	\end{aligned}
	\]
	$F_k(x)=0$ if $x<p_k$.

	\textbf{Prime sieve.}
	Define $G_{d,k}(x)=\sum m^d$ over $2\le m\le x$ where
	$m$ is prime or $\operatorname{lpf}(m)>p_k$.
	Initially $G_{d,0}(x)=\sum_{m=2}^x m^d$.
	\[
	\begin{aligned}
	G_{d,k}(x)={}&G_{d,k-1}(x)\\
	&-[p_k^2\le x]p_k^d
	\left(G_{d,k-1}(\lfloor x/p_k\rfloor)
	-\sum_{i<k}p_i^d\right).
	\end{aligned}
	\]
	After sieving $p\le\sqrt n$,
	$G_d(x)=\sum_{p\le x,\ p\text{ prime}}p^d$.
	If $f(p)=\sum_d a_dp^d$, then
	$F_{\mathrm{prime}}(x)=\sum_d a_dG_d(x)$.

	\textbf{Optional DP.}
	Define $H_k(x)=1+F_k(x)$.
	\[
	H_k(x)=H_{k+1}(x)+
	\sum_{\substack{c\ge1\\p_k^c\le x}}
	f(p_k^c)H_{k+1}(\lfloor x/p_k^c\rfloor).
	\]
	Let $r=\pi(\sqrt n)$.
	Initialize $H_{r+1}(x)=1+F_{\mathrm{prime}}(x)
	-F_{\mathrm{prime}}(\min(x,p_r))$.

	\textbf{Answer.}
	$S(n)=\sum_{i=1}^n f(i)=1+F_1(n)=H_1(n)$.
All divisions are integer divisions.

	\kactlimport{Min25.h}

\section{Divisibility}
	\kactlimport{euclid.h}
	\kactlimport{Euclid.java}
	\kactlimport{CRT.h}

	\subsection{Bézout's identity}
	For $a \neq $, $b \neq 0$, then $d=gcd(a,b)$ is the smallest positive integer for which there are integer solutions to
	$$ax+by=d$$
	If $(x,y)$ is one solution, then all solutions are given by
	$$\left(x+\frac{kb}{\gcd(a,b)}, y-\frac{ka}{\gcd(a,b)}\right), \quad k\in\mathbb{Z}$$

	\kactlimport{phiFunction.h}

\section{Fractions}
	\kactlimport{ContinuedFractions.h}
	\kactlimport{FracBinarySearch.h}

\section{Pythagorean Triples}
 The Pythagorean triples are uniquely generated by
 \[ a=k\cdot (m^{2}-n^{2}),\ \,b=k\cdot (2mn),\ \,c=k\cdot (m^{2}+n^{2}), \]
 with $m > n > 0$, $k > 0$, $m \bot n$, and either $m$ or $n$ even.

\section{Primes}
	$p=962592769$ is such that $2^{21} \mid p-1$, which may be useful. For hashing
	use 970592641 (31-bit number), 31443539979727 (45-bit), 3006703054056749
	(52-bit). There are 78498 primes less than 1\,000\,000.

	Primitive roots exist modulo any prime power $p^a$, except for $p = 2, a > 2$, and there are $\phi(\phi(p^a))$ many.
	For $p = 2, a > 2$, the group $\mathbb Z_{2^a}^\times$ is instead isomorphic to $\mathbb Z_2 \times \mathbb Z_{2^{a-2}}$.

\section{Estimates}
	$\sum_{d|n} d = O(n \log \log n)$.

	The number of divisors of $n$ is at most around 100 for $n < 5e4$, 500 for $n < 1e7$, 2000 for $n < 1e10$, 200\,000 for $n < 1e19$.

\section{Mobius Function}
\[
	\mu(n) = \begin{cases} 0 & n \textrm{ is not square free}\\ 1 & n \textrm{ has even number of prime factors}\\ -1 & n \textrm{ has odd number of prime factors}\\\end{cases}
\]
  Mobius Inversion:
  \[ g(n) = \sum_{d|n} f(d) \Leftrightarrow f(n) = \sum_{d|n} \mu(d)g(n/d) \]
  Other useful formulas/forms:

  $ \sum_{d | n} \mu(d) = [ n = 1] $ (very useful)

  $ g(n) = \sum_{n|d} f(d) \Leftrightarrow f(n) = \sum_{n|d} \mu(d/n)g(d)$

 $ g(n) = \sum_{1 \leq m \leq n} f(\left\lfloor\frac{n}{m}\right \rfloor ) \Leftrightarrow f(n) = \sum_{1\leq m\leq n} \mu(m)g(\left\lfloor\frac{n}{m}\right\rfloor)$
